\chapter{Linking, Loading, ABI, and Whole-Program Optimisation}

The compiler sees one translation unit at a time; the linker assembles them, the loader maps the program at run time, and the ABI is the binary contract that lets independently-compiled pieces interoperate. This layer is invisible until an undefined-reference error, a silent ODR violation, a slow startup, or a failed cross-TU inline reveals it. This chapter explains the linking and loading model, symbol visibility and ABI stability, and the whole-program optimisations (LTO, PGO) that recover cross-TU performance.

\section*{Chapter Roadmap}
\begin{itemize}
\item 102.1 The Compilation and Linking Model
\item 102.2 Symbol Visibility and Name Mangling
\item 102.3 Static vs Dynamic Linking and the Loader
\item 102.4 PIC, PIE, and Relocations
\item 102.5 The ABI and Its Stability
\item 102.6 Whole-Program Optimisation: LTO and PGO
\item 102.7 Build Hazards and the Discipline
\end{itemize}

\section{The Compilation and Linking Model}
Each TU compiles independently to an object file with defined and undefined symbols; the linker resolves every undefined symbol, merges sections, and produces the binary.

\begin{lstlisting}[language=bash, caption={Separate compilation then linking.}]
g++ -c a.cpp -o a.o
g++ -c b.cpp -o b.o
g++ a.o b.o -o app
\end{lstlisting}

\textbf{Why this matters.} \texttt{undefined reference to X} means a used symbol was never defined; \texttt{multiple definition of X} means two TUs defined a non-inline symbol (an ODR violation, Chapter 104). The per-TU view is also why the optimizer cannot inline across TUs by default (102.6).

\section{Symbol Visibility and Name Mangling}
C++ mangles names (encoding the signature) so overloads get distinct symbols; \texttt{extern "C"} disables mangling for C interop.

\begin{lstlisting}[language=C++, caption={Mangling distinguishes overloads; extern "C" gives a plain symbol. Min standard: C++98.}]
void f(int);                  // _Z1fi
void f(double);               // _Z1fd
extern "C" void c_api(int);   // symbol "c_api", no mangling/overloading
\end{lstlisting}

Visibility controls which symbols a shared library exports; \texttt{-fvisibility=hidden} keeps internal symbols private.

\textbf{Why this matters / cost model.} Mismatched declarations fail to link rather than misbehave. Exporting every symbol bloats the dynamic symbol table (slow load/relocation), risks ODR clashes, and blocks inlining of exported functions. Hiding internals shrinks the table, speeds loading, and unlocks optimization.

\section{Static vs Dynamic Linking and the Loader}
Static linking copies code into the executable; dynamic linking records a dependency and the loader maps the \texttt{.so} at run time, resolving symbols.

\begin{itemize}
\item Static: larger binary, fast startup, duplicated memory, self-contained, recompile to update.
\item Dynamic: smaller binary, slower startup (load+relocate+resolve), shared memory, needs the libs, swap to update.
\end{itemize}

\textbf{Why this matters / cost model.} The loader maps each \texttt{.so}, applies relocations, and resolves symbols (lazily via the PLT or eagerly with \texttt{-z now}) --- measurable startup latency and a tiny per-call PLT cost. Static linking eliminates all of it and lets LTO see the library code, hence its use in latency-critical binaries. Trade-offs: size, no shared memory, no hot library patching.

\section{PIC, PIE, and Relocations}
Shared libraries and PIE executables use position-independent code, accessing globals via the GOT and exported functions via the PLT; the loader applies relocations to fix these up.

\textbf{Why this matters / cost model.} PIC/PIE add a GOT indirection per global and a PLT jump per inter-module call (lazily resolved) --- negligible mostly, but a hot loop touching PI globals can feel it (\texttt{-fno-plt}, static linking help). The larger cost is load-time relocation for binaries with many symbols/libraries. Reduce exported symbols, prefer static linking, and choose \texttt{-z now} vs lazy binding. PIE enables ASLR (a security feature).

\section{The ABI and Its Stability}
The ABI is the binary contract: calling convention, name mangling, class layout, vtable layout, exception propagation. Components interoperate only if they agree.

\textbf{Why this matters.} ABI stability lets you link libraries built by other compiler versions or load plugins without recompiling --- but it is a cage: the standard library ABI is frozen because changing \texttt{std::string}'s layout would break every old binary. Practical consequences: never mix incompatible ABIs (\texttt{\_GLIBCXX\_USE\_CXX11\_ABI} mismatch corrupts silently); be cautious passing standard types across uncontrolled \texttt{.so} boundaries; and design your own ABIs with Pimpl/pure-virtual interfaces and \texttt{extern "C"} factories. An ABI break is a recompile-the-world event.

\section{Whole-Program Optimisation: LTO and PGO}
LTO emits IR into objects and re-runs the optimizer at link time with the whole program visible, inlining across TUs. PGO feeds a profiling run back into codegen.

\begin{lstlisting}[language=bash, caption={LTO enables cross-TU inlining; PGO feeds real execution data into codegen.}]
g++ -O2 -flto -c a.cpp b.cpp
g++ -O2 -flto a.o b.o -o app
g++ -O2 -fprofile-generate app.cpp -o app && ./app <representative-input>
g++ -O2 -fprofile-use     app.cpp -o app
\end{lstlisting}

\textbf{Why this matters / cost model.} LTO is how zero-cost abstraction survives separate compilation: a hot cross-TU accessor that could not inline now does, often a 5--15\% win plus dead-code reductions. PGO lays out the binary for I-cache locality and inlines genuinely-hot callees. Costs: LTO increases link time and memory; PGO needs a representative workload (a bad profile can pessimize). Enable LTO broadly; add PGO for the hottest services.

\section{Build Hazards and the Discipline}
\begin{itemize}
\item \texttt{undefined reference} --- link the library; match declarations.
\item \texttt{multiple definition} --- \texttt{inline}/\texttt{static}; one definition (Chapter 104).
\item Silent ABI mismatch --- consistent toolchain/flags; stable public ABI.
\item Slow startup --- \texttt{-fvisibility=hidden}; static link; \texttt{-z now}.
\item Abstraction not free --- enable LTO.
\item Wrong hot-path layout --- PGO with representative input.
\end{itemize}

\textbf{The discipline.} Mastering this layer pays off in correctness (mangling, ODR, ABI), startup latency (visibility, static linking), and steady-state performance (LTO, PGO). Defaults for a performance-critical binary: hide internal symbols, prefer static linking, enable LTO in release, add PGO for the hottest paths, and treat your public ABI as a long-lived contract.
